\section{Introduction and Career Interest}

This report studies the Site Reliability Engineer (SRE) position at Google as a target for cooperative education (COOP) or job training. It answers the assignment questions in order: what the job involves, which skills and tools it needs, which operating systems, hardware and infrastructure it works with, and how this connects to the Operating Systems course.

\paragraph{Why SRE.}
Our reason is our own judgment about where computing work is going. It is an opinion, not a verified fact. Today's Generative AI tools help most with writing software. Keeping that software running in production, on real machines and networks, still needs engineers who can tell what is wrong with a live system. We expect this kind of work to stay valuable as Generative AI spreads. We have not tested this expectation against labor-market data, so we treat it as an assumption. The role also fits the course. Google's SRE book says that SRE is what happens when a software engineer is asked to design an operations team \cite{sre-book}. The job therefore joins programming with operating systems, networking and hardware.

\paragraph{Why Google.}
We chose Google for two reasons. First, it is a large company with a long SRE history. The SRE book's introduction says its author joined Google in 2003 and was tasked with running a ``Production Team'' of seven engineers \cite{sre-book}. Second, Google publishes its SRE practice openly at \texttt{sre.google}: the SRE book, the SRE Workbook, \emph{Building Secure and Reliable Systems} and further articles. This let us check our claims against Google's own words and not only against second-hand summaries.

\paragraph{Limits of this choice.}
The same sources also limit what we can say. The books are Google's own publications, dated 2017 to 2020, so most of what we report about Google comes from one company's account. They describe Google's practice when they were written, and they say little about the tools Google uses today. We therefore state a claim about Google only where a published source supports it. Where no published source describes Google's internal environment, we describe a typical environment for the position and label it as an assumption or recommendation. The one Google SRE internship that we found is for the EMEA region only (Section~\ref{sec:org}), so we treat Google as a study target.
